\clearpage

\section{Spectral Mass Spectrum and Hierarchy}
  \label{sec:spectral-mass-hierarchy}

  Within the spectral admissibility framework established in
  Chapter~~\ref{sec:spectral-admissibility}, particle masses appear as invariants associated with
  stable modes of the projection fiber.
  The mass hierarchy therefore reflects the structure of admissible
  spectral configurations rather than independent dynamical
  parameters.

  This section develops the physical consequences of that structure.
  We first interpret rest mass as a spectral stability condition of
  projected configurations.
  We then show how torsional constraints generate mass amplification,
  how higher torsion sectors may lead to irreducible spectral
  structures, and how related
  mechanisms may leave large-scale signatures in the form of
  gravitational spectral wakes.
  Throughout this chapter, torsion sectors are labelled by an abstract
  index~$r$.
  This index is distinct from the electromagnetic fiber winding~$w$ of
  Section~\ref{sec:topological-constraints}, and no identification of a
  torsion sector with a specific particle is made in this chapter.

  Finally, the same spectral framework provides a geometric
  interpretation of flavor structure and CP violation as arising
  from incompatible spectral extractions within the same projection
  fiber.

  \subsection{Spectral Stability and the Unit of Mass}
  \label{sec:spectral-stability}

  Within the admissibility framework developed in Chapter~\ref{sec:spectral-admissibility},
  rest mass appears as a spectral invariant associated with
  the stability of projected configurations.

  In the effective description, particle states correspond to
  stable modes supported by the projection fiber~$\Pi$.
  Their mass is determined by the eigenvalues of the scalar
  Laplacian~$\Delta^{(0)}_G$ acting on this space, subject to
  topological constraints~$\mathcal{T}$:
  \begin{equation}
    m^2 c^2 \;=\; \lambda_{\mathcal{T}}
    \;\equiv\;
    \mathrm{Eig}\!\left(\Delta^{(0)}_G\right)
    \Big|_{\mathcal{T}} .
    \label{eq:mass_spectral_def}
  \end{equation}

  The quantity $\lambda_{\mathcal{T}}$ represents the minimal
  spectral cost required to maintain the corresponding
  configuration within the admissible projective sector.
  Mass therefore measures the resistance of a configuration
  to projective ordering advancement.

  The reference mass scale~$m_{(1)}$ corresponds to the lowest
  non-trivial admissible eigenmode~$\lambda_1$ of the fiber
  geometry $\Pi \cong S^3$.
  All other sector mass scales can then be expressed relative to
  this fundamental scale,
  \begin{equation}
    m \;=\; m_{(1)}
    \sqrt{\frac{\lambda_{\mathcal{T}}}{\lambda_1}} .
  \end{equation}

  \subsection{Non-Commutativity as a Source of Mass}
  \label{sec:noncommutativity-mass}

  Within the spectral framework established in Chapter~\ref{sec:spectral-admissibility},
  torsion acts as a dynamical constraint that competes with
  relaxation transport through the projection fiber.

  \subsubsection*{Inhibition of Projective Ordering}
    \label{sec:inhibition-relaxation}

    Let $\Omega_r$ denote the effective torsion operator associated
    with the abstract torsion sector index~$r$.
    For the fundamental sector ($r=1$),
    torsion and projective ordering remain spectrally compatible:
    \begin{equation}
      [\Delta^{(0)}_G,\Omega_{1}] = 0 .
    \end{equation}

    For higher torsion sectors ($r\ge 2$), torsional constraints
    become spectrally frustrated:
    \begin{equation}
      [\Delta^{(0)}_G,\Omega_{r}] \neq 0
      \qquad (r\ge 2).
    \end{equation}

    This non-commutativity inhibits uniform projective ordering across the
    projection fiber and induces an effective spectral compression,
    which manifests as increased inertial mass.

    The torsional contribution can be expressed through the spectral
    invariant
    \begin{equation}
      \mathcal{A}(r) \equiv
      \frac12 \ln\!\left(
                     \frac{\det_\zeta(\Delta^{(0)}_G+\Omega_{r})}
                     {\det_\zeta(\Delta^{(0)}_G)}
      \right),
      \label{eq:Aw-fredholm}
    \end{equation}
    where $\det_\zeta$ denotes zeta-regularized determinants.

  \subsubsection*{Dynamical Stability of Torsional Frequencies}
    \label{sec:pisano-ratio}

    For $r=2$, the projection fiber becomes dynamically anisotropic
    and splits into competing sectors
    \[
      \Pi=\Pi_{\parallel}\oplus\Pi_{\perp}.
    \]

    The quantities $\lambda_{\parallel}$ and $\lambda_{\perp}$ are
    effective torsional rotation frequencies generated by $\Omega_2$
    within the first multiplet.
    They are extracted from the commutator flow on observables
    restricted to $\mathcal{H}_{k=1}$.

    Dynamical stability favors frequency ratios that remain robust
    under perturbations.
    In analogy with KAM-type criteria, the most stable admissible
    ratio corresponds to the golden ratio
    \begin{equation}
      \frac{\lambda_{\parallel}}{\lambda_{\perp}}
      = \varphi ,
      \qquad
      \varphi=\frac{1+\sqrt5}{2}.
      \label{eq:beta-golden}
    \end{equation}

    The appearance of $\varphi$ therefore reflects a stability
    selection rule rather than a static spectral extremum.

  \subsubsection*{Low-Sector Spectrum Synthesis}
    \label{sec:leptonic-synthesis}

    The low-sector mass hierarchy combines two distinct mechanisms:

    \begin{itemize}
      \item a topologically determined spectral invariant
      \[
        \frac{\lambda_2}{\lambda_1}=\frac{8}{3},
      \]

      \item a dynamical stability selection rule governed by
      torsional non-commutativity.
    \end{itemize}

    For the $r=2$ sector this leads to the parameter-free relation
    \begin{equation}
      \boxed{
        \frac{m_{(2)}}{m_{(1)}}
        = \sqrt{\frac{\lambda_{2}}{\lambda_{1}}}
        \cdot \frac{3}{2\alpha}
        \cdot \frac{1}{\varphi}
      }
      \qquad\text{with}\qquad
      \frac{\lambda_2}{\lambda_1}=\frac{8}{3}.
      \label{eq:muon-ratio-final}
    \end{equation}

    The robustness of $\lambda_2/\lambda_1=8/3$ is demonstrated
    independently in Appendix~\ref{sec:spectral_ratio_derivation}.

    \paragraph*{Outlook.}
      The same commutator-driven torsional mechanism extends naturally
      to higher torsion sectors ($r\ge 3$), providing a route to the
      mass scales of composite sectors.

  \subsection{Irreducible Torsion and Higher Torsion Sectors}
  \label{sec:irreducible-torsion}

  Within the spectral framework established in Chapter~\ref{sec:spectral-admissibility},
  the condition
  \(
  [\Delta^{(0)}_G,\Omega_r]\neq 0
  \)
  signals spectral frustration for all $r\ge 2$.
  However, non-commutativity alone does not distinguish the qualitative
  structure of different higher torsion sectors.
  A further structural criterion is required.

  \subsubsection*{From Anisotropy to Irreducibility}

    In the $r=2$ sector, torsion induces anisotropy within the projection
    fiber, leading to a controlled splitting
    \(
    \Pi = \Pi_{\parallel}\oplus\Pi_{\perp}.
    \)
    The associated dynamical selection principle therefore operates
    within a reducible decomposition of the first spectral multiplet.

    For $r\ge 3$, the situation may be qualitatively different.
    The torsional operator $\Omega_r$ can act on the admissible subspace
    in a manner that is not decomposable into independent lower-index
    sectors.

    Formally, let
    \begin{equation}
      \mathcal{A}_r
      \equiv
      \mathrm{alg}\!\left(
                      \Delta^{(0)}_G,\,\Omega_r
      \right)
      \Big|_{H_{\mathrm{adm}}^{(r)}}
    \end{equation}
    denote the algebra generated by the projective ordering operator and torsion,
    restricted to the admissible configuration space.

    If the commutant
    \begin{equation}
      \mathcal{A}_r'
      =
      \big\{
      X\in\mathrm{End}(H_{\mathrm{adm}}^{(r)})
      \mid
      [X,A]=0\ \forall A\in\mathcal{A}_r
      \big\}
    \end{equation}
    reduces to scalar multiples of the identity,
    the torsional action is irreducible in the admissible sector.

    Irreducibility implies that no non-trivial projector exists that
    simultaneously preserves admissibility and diagonalizes the torsional
    action.
    In that case a naive decomposition
    \(
    \Omega_r \simeq \Omega_{1}^{\oplus r}
    \)
    is not dynamically accessible within $H_{\mathrm{adm}}^{(r)}$.

    Operationally, irreducibility can be tested by computing the
    dimension of the commutant $\mathcal{A}_r'$.
    If $\dim \mathcal{A}_r' = 1$, the torsional sector is algebraically
    irreducible within $H_{\mathrm{adm}}^{(r)}$.

  \subsubsection*{Admissibility as a Structural Constraint}

    The admissible subspace $H_{\mathrm{adm}}^{(r)}$ is defined by the
    bounded-admissibility condition.
    Factorizations that would decompose $\Omega_r$ into independent
    lower-index components may require intermediate configurations
    that violate admissibility.

    In such cases irreducibility is not an additional axiom but a
    consequence of bounded admissibility:
    the decomposition path exits the admissible domain and is therefore
    dynamically forbidden.

    This mechanism provides a structural explanation for why higher
    torsion sectors may exhibit stronger mass amplification than the
    $r=2$ case.
    The obstruction is therefore not merely energetic but algebraic.

  \subsubsection*{Spectral Consequences for $r=3$}

    The $r=3$ sector is the first case in which genuine algebraic
    irreducibility may arise.
    If $\mathcal{A}_3'$ is trivial within $H_{\mathrm{adm}}^{(3)}$,
    the torsional action cannot be reduced to a sum of independent
    $r=1$ sectors.

    In that regime the relevant mass scale is controlled by a global
    spectral invariant rather than by a two-mode dynamical selection rule.
    A natural candidate is the torsional action
    \(
    \mathcal{A}(3)
    \)
    defined in Eq.~\eqref{eq:Aw-fredholm}, or an associated canonical
    ratio extracted from the spectrum of
    \(
    \Delta^{(0)}_G+\Omega_3
    \)
    restricted to $H_{\mathrm{adm}}^{(3)}$.

    Identifying such an invariant constitutes the next structural step
    beyond the reducible sectors.

  \subsection{The Irreducible Sector and Spectral Hierarchy}
  \label{sec:baryonic-sector}

  Within the spectral admissibility framework established in Chapter~\ref{sec:spectral-admissibility},
  the $r=3$ configuration is the first sector in which algebraic
  irreducibility becomes structurally possible within the admissible
  subspace.
  If no non-trivial projector simultaneously preserves admissibility and
  diagonalizes $\mathcal{A}_3$, the resulting mass scale cannot be
  interpreted as a superposition of lower-sector contributions.

  \subsubsection*{From Dynamical Selection to Global Invariants}

    In the $r=2$ sector, mass amplification arises from a dynamical
    selection rule acting within a reducible two-mode structure.
    The Pisano ratio $\varphi$ enters as a stability condition on the
    torsional generator restricted to the first multiplet.

    For $r=3$, such a two-mode reduction need not exist.
    If the torsional action is algebraically irreducible in
    $H_{\mathrm{adm}}^{(3)}$, no analogue of the Pisano selection rule
    applies, because no non-trivial invariant subspace remains on which
    a two-mode frequency selection could operate.
    The relevant observable must then be a global spectral invariant.

    A natural candidate is the torsional action
    \(
    \mathcal{A}(3)
    \)
    defined in Eq.~\eqref{eq:Aw-fredholm},
    or an invariant ratio extracted canonically from the spectrum of
    \(
    \Delta^{(0)}_G+\Omega_3
    \)
    restricted to $H_{\mathrm{adm}}^{(3)}$.
    Such an invariant should be independent of overall amplitude
    rescalings of $\Omega_3$ and stable under increasing spectral
    truncation.

  \subsubsection*{Mass Hierarchy as Spectral Scaling}

    If the $r=3$ sector is irreducible,
    the sector mass scale is controlled by the spectral distortion
    induced by $\Omega_3$.
    In analogy with the reducible sectors one expects
    \begin{equation}
      \frac{m_{(3)}}{m_{(1)}}
      \;\sim\;
      \sqrt{
        \frac{\lambda_{\mathrm{bind}}^{(3)}}
        {\lambda_{1}}
      }.
    \end{equation}

    Unlike the reducible case,
    this ratio would not be determined by a two-mode dynamical fixed point,
    but by the global spectral structure of the irreducible torsional
    sector.

  \subsubsection*{Programmatic Outlook}

    The structural task is therefore twofold:

    \begin{itemize}
      \item[(i)] determine whether $\mathcal{A}_3'$ is trivial within
      $H_{\mathrm{adm}}^{(3)}$,
      \item[(ii)] identify a canonical spectral invariant of
      $\Delta^{(0)}_G+\Omega_3$ that is stable under admissible
      deformations.
    \end{itemize}

    If such an invariant exists and does not depend on parameter tuning,
    the irreducible-sector mass scale would emerge as a structural
    consequence of spectral irreducibility rather than as
    phenomenological input.

  \input{5-predictions_discussion_conclusion/11-spectral-mass-hierarchy/005-gravitational-shadows}
  \input{5-predictions_discussion_conclusion/11-spectral-mass-hierarchy/006-flavor-structure-cp}
